\documentclass[11pt]{article}
\usepackage[margin=1in]{geometry}
\usepackage[T1]{fontenc}
\usepackage{graphicx}
\usepackage{hyperref}
\usepackage{enumitem}
\setlength{\parskip}{0.5em}
\setlength{\parindent}{0pt}

\title{graph neural networks}
\author{arXiv GraphRAG --- phase-6 deep research}
\date{\today}

\begin{document}
\maketitle

\begin{abstract}
The review is excellent in its narrative flow and LLM integration but misses key recent sub-fields like equivariant GNNs and temporal graphs.
\end{abstract}

\section{The Evolution of Relational Intelligence: A Comprehensive Review of Graph Neural Networks}

The representation of complex relational data has long posed a significant challenge to classical machine learning paradigms, which are primarily optimized for Euclidean structures such as grids or sequences. As real-world datasets—ranging from molecular structures to social networks and knowledge graphs—increasingly manifest as irregular, non-Euclidean topologies, the necessity for models capable of capturing intricate relational dependencies has become paramount. Graph Neural Networks (GNNs) have emerged as the definitive solution to this problem, providing a mathematical framework to learn sophisticated representations directly from graph-structured data. The ability to effectively model these dependencies is critical, as it underpins modern breakthroughs in fields as diverse as drug discovery, recommendation systems, and social computing.

This review traces the historical and technical trajectory of the field, beginning with the foundational principles of early embeddings and the subsequent architectural evolution from Graph Convolutional Networks (GCNs) to generalized message-passing frameworks. We examine the theoretical constraints that define the field's capabilities, specifically focusing on the limits of expressivity, stability, and structural learning. The discussion then expands into the domain of geometric deep learning, exploring how equivariance and symmetry are leveraged to handle complex spatial relationships. Furthermore, we analyze the shift toward modern training paradigms, including self-supervised, contrastive, and generative learning, alongside the critical necessity for explainability and interpretability in graph-based decision-making. Finally, we look toward the current frontier, where the integration of graph transformers and large language models (LLMs) is redefining the boundaries of relational intelligence.

\subsection{Foundations and Early Embeddings}

Early research in graph representation learning can be categorized into two distinct paradigms: random-walk-based embedding methods and spectral-based convolutional architectures.

The embedding paradigm began with DeepWalk [\texttt{arXiv:1403.6652}], which treats truncated random walks as sentences for language modeling. This approach proved particularly robust in sparse data scenarios, achieving up to 10\% higher F1 scores while using 60\% less training data than baselines [\texttt{arXiv:1403.6652}]. To address scalability, LINE [\texttt{arXiv:1503.03578}] introduced an edge-sampling algorithm to optimize local and global structure preservation, enabling the method to scale to millions of nodes and billions of edges [\texttt{arXiv:1503.03578}]. Building upon these foundations, node2vec [\texttt{arXiv:1607.00653}] generalized the random walk approach by introducing a flexible biased procedure, allowing for more expressive, task-independent node representations [\texttt{arXiv:1607.00653}].

In contrast, the spectral paradigm focuses on applying convolutional operations directly to graph structures. Early work established the use of the graph Laplacian spectrum and hierarchical clustering to create convolutional layers with parameter counts independent of input size [\texttt{arXiv:1312.6203}]. This was significantly advanced by the Graph Convolutional Network (GCN) [\texttt{arXiv:1609.02907}], which utilizes a localized first-order approximation of spectral convolutions and a renormalization trick for numerical stability. While GCN achieves superior classification performance—notably 81.5\% accuracy on the Cora dataset [\texttt{arXiv:1609.02907}]—it faces practical limitations, such as scalability constraints due to full-batch gradient descent memory requirements and a lack of native support for directed edges or edge features [\texttt{arXiv:1609.02907}]. Subsequent developments like Graph Attention Networks (GAT) [\texttt{arXiv:1710.10903}] extended the GCN framework by incorporating masked self-attentional layers, allowing nodes to weight their neighbors differently without requiring prior knowledge of the graph structure or costly matrix operations [\texttt{arXiv:1710.10903}].

A critical shift occurred when moving from transductive to inductive learning. While earlier embedding methods like DeepWalk [\texttt{arXiv:1403.6652}] and node2vec [\texttt{arXiv:1607.00653}] are primarily transductive, GraphSAGE [\texttt{arXiv:1706.02216}] introduced an inductive framework that learns aggregator functions to generate embeddings for previously unseen nodes via neighborhood sampling [\texttt{arXiv:1706.02216}]. This approach is approximately 100× faster than transductive baselines at inference and improves supervised F1-scores by an average of 51\% over raw features [\texttt{arXiv:1706.02216}]. However, the theoretical foundation of GraphSAGE is somewhat limited, as its proof for learning structural information relies on the restrictive assumption that all node features are distinct [\texttt{arXiv:1706.02216}].

\subsection{Architectural Evolution: From GCNs to Message Passing}

The architectural landscape of Graph Neural Networks (GNNs) has transitioned from rigid spectral frameworks toward flexible, message-passing mechanisms capable of handling diverse relational complexities. Early developments focused on extending spectral networks to non-Euclidean domains by incorporating graph estimation procedures to reduce parameter counts and learning complexity [\texttt{arXiv:1506.05163}]. This period established the theoretical foundation for geometric deep learning, framing the challenge of processing data on graphs and manifolds [\texttt{arXiv:1611.08097}].

A significant shift occurred with the introduction of gated mechanisms for sequence modeling. While spectral methods often focused on graph-level tasks [\texttt{arXiv:1506.05163}], Gated Graph Sequence Neural Networks introduced Gated Recurrent Units (GRUs) to facilitate message passing, allowing for the modeling of sequences of graph-structured inputs [\texttt{arXiv:1511.05493}]. This move toward gated architectures provided the flexibility needed for complex applications like program verification and graph algorithm learning [\texttt{arXiv:1511.05493}].

As the field matured, architectures evolved to address specific data structures, such as multi-relational knowledge bases. Relational Graph Convolutional Networks (R-GCNs) moved beyond simple connectivity by designing message passing specifically for multi-relational data, yielding a 29.8\% improvement on the FB15k-237 dataset for link prediction [\texttt{arXiv:1703.06103}]. However, standard GNNs face fundamental limitations in expressive power. While some models attempt to expand this capacity by utilizing \$k\$-dimensional Weisfeiler-Leman (k-WL) heuristics to capture higher-order structures [\texttt{arXiv:1810.02244}], others have identified critical failures in depth. Specifically, it has been demonstrated that GCNs can exponentially lose expressive power for node classification as layers increase, a phenomenon tied to the spectra of the normalized Laplacian [\texttt{arXiv:1905.10947}].

To mitigate these limitations, researchers have explored hybrid propagation strategies. For instance, combining GCNs with Personalized PageRank (PPR) via the PPNP framework allows for the use of large, adjustable neighborhoods, providing a more robust mechanism for semi-supervised classification [\texttt{arXiv:1810.05997}]. Furthermore, the characterization of GNNs through Graph Signal Processing (GSP) has enabled the design of graph convolutional filters that maintain permutation equivariance and stability against topological changes [\texttt{arXiv:2003.03777}]. This evolution highlights a trajectory from basic spectral filters toward sophisticated, higher-order, and signal-aware message-passing architectures that balance expressive power with structural stability.

\subsection{Expressivity, Stability, and Structural Learning}

The evolution of Graph Neural Network (GNN) research has been fundamentally driven by the quest to characterize and expand their expressive power relative to classical graph isomorphism tests. A seminal contribution to this field is the establishment of a theoretical framework that links the expressive power of neighborhood aggregation schemes to the Weisfeiler-Lehman (WL) graph isomorphism test [\texttt{arXiv:1810.00826}]. By utilizing recursive aggregation and transformation of neighboring node representations, this work demonstrates that standard GNNs are provably as powerful as the 1-WL test [\texttt{arXiv:1810.00826}].

Subsequent research has sought to transcend these 1-WL limits by developing higher-order architectures. For instance, k-order invariant and equivariant networks can match k-WL expressiveness by interleaving Multilayer-Perceptron (MLP) applications with quadratic matrix multiplication operations [\texttt{arXiv:1905.11136}]. Specifically, a scalable 2-order network was proposed that guarantees 3-WL power [\texttt{arXiv:1905.11136}]. This pursuit of higher-order expressivity is further refined by the proof of equivalence between graph isomorphism testing and permutation-invariant function approximation, which introduced the Ring-GNN to distinguish non-isomorphic regular graphs [\texttt{arXiv:1905.12560}].

The limitations of these models are often analyzed through the lens of substructure counting. While standard Message Passing Neural Networks (MPNNs) are constrained by the expressive limits of 2-WL and 2-invariant graph networks (2-IGNs), the introduction of Local Relational Pooling has helped establish these boundaries [\texttt{arXiv:2002.04025}]. To move beyond these constraints, Graph Substructure Networks (GSN) utilize a topologically-aware message passing scheme based on substructure encoding, achieving expressivity strictly greater than the WL test while maintaining linear network complexity and locality [\texttt{arXiv:2006.09252}].

However, increasing expressivity does not inherently resolve issues regarding information flow and stability. The phenomenon of "over-squashing" occurs when graph bottlenecks impede effective message passing [\texttt{arXiv:2111.14522}]. This can be analyzed through edge-based combinatorial curvature, which provides a precise description of how bottlenecks affect node feature propagation [\texttt{arXiv:2111.14522}]. To mitigate these stability and flow issues, research has pivoted toward hybrid architectures. For example, GraphTrans addresses the need for long-range context by appending a permutation-invariant Transformer module to a standard GNN, allowing the model to learn long-range pairwise relationships without explicit structural encoding [\texttt{arXiv:2201.08821}]. Thus, the literature reveals a tension between local structural expressivity—achieved through higher-order WL-equivalent mechanisms [\texttt{arXiv:1905.11136}, \texttt{arXiv:2002.04025}]—and the necessity of global context to overcome topological bottlenecks [\texttt{arXiv:2111.14522}, \texttt{arXiv:2201.08821}].

\subsection{Geometric Deep Learning and Equivariance}

The evolution of graph neural networks (GNNs) has transitioned from general structural modeling to the rigorous integration of geometric symmetries. Early work in graph learning often focused on structural properties, such as the deep autoregressive modeling of node and edge formations in GraphRNN (\texttt{arXiv:1802.08773}) or the use of score functions like the Katz index for link prediction (\texttt{arXiv:1802.09691}). However, as applications shifted toward scientific domains like molecular modeling and dynamical systems, the necessity for formalizing invariance and equivariance became paramount.

A foundational theoretical contribution is provided by \texttt{arXiv:1812.09902}, which characterizes all permutation invariant and equivariant linear layers for (hyper-)graph data. By utilizing Bell numbers for dimension characterization and providing orthogonal bases, this work proves that such models can approximate any message passing neural network. While this establishes a framework for permutation equivariance, subsequent research has expanded these principles to include Euclidean symmetries. For instance, \texttt{arXiv:2102.09844} introduces E(n) Equivariant Graph Neural Networks (EGNNs), which achieve equivariance to rotations, translations, reflections, and permutations. A significant strength of EGNNs is their ability to maintain scalability to higher-dimensional spaces without the computational burden of higher-order representations, making them highly effective for predicting molecular properties.

Recent advancements have sought to integrate these geometric principles into more complex architectures. The Equiformer (\texttt{arXiv:2206.11990}) represents a sophisticated leap by replacing standard dot-product attention with an equivariant graph attention mechanism. Unlike the more streamlined EGNNs, Equiformer utilizes non-linear message passing through tensor products and relies on SE(3)/E(3)-equivariance based on irreducible representations (irreps) to model 3D atomistic graphs. This allows for a more nuanced representation of 3D structures compared to models that focus solely on E(n) symmetries.

The complexity of this field is reflected in recent systematic reviews. \texttt{arXiv:2202.07230} provides a survey that classifies geometrically equivariant GNNs into three distinct groups based on their message passing and aggregation mechanisms, specifically focusing on scientific problems involving geometric graphs. This is complemented by the more recent survey in \texttt{arXiv:2403.00485}, which offers a broader formalization of geometric graphs. While \texttt{arXiv:2202.07230} emphasizes the classification of existing methods, \texttt{arXiv:2403.00485} provides a unified view of models through the lens of geometric message passing, encompassing physical symmetries—such as translations, rotations, and reflections—and the underlying geometry and topology of the data. Together, these works trace a trajectory from simple permutation invariance to the complex, symmetry-aware architectures required for modern scientific machine learning.

\subsection{Self-Supervised, Contrastive, and Generative Learning}

The landscape of graph representation learning has expanded from supervised methods to include diverse self-supervised learning (SSL) paradigms, which can be categorized into generation-based, auxiliary property-based, contrast-based, and hybrid approaches [\texttt{arXiv:2103.00111}].

Generative models aim to capture the underlying distribution of graph structures and attributes. Early efforts focused on Variational Autoencoders (VAEs), such as GraphVAE, which sidesteps linearization hurdles by directly outputting a probabilistic fully-connected graph of a predefined maximum size [\texttt{arXiv:1802.03480}]. This approach is specialized for molecule generation [\texttt{arXiv:1802.03480}]. Building upon these concepts, more general approaches have emerged to learn generative models over arbitrary graphs, utilizing graph neural networks to express probabilistic dependencies among both nodes and edges [\texttt{arXiv:1803.03324}]. While these VAE-based methods provide a foundation for capturing structure and attributes across domains like knowledge graphs and chemistry [\texttt{arXiv:1803.03324}], recent advancements have shifted toward diffusion-based frameworks. For instance, DiGress utilizes a discrete denoising diffusion process and a Markovian noise model to handle graphs with categorical attributes [\texttt{arXiv:2209.14734}]. Unlike earlier generative models, DiGress scales effectively to large datasets, such as 1.3M molecules, by preserving marginal distributions through a Graph Transformer Network architecture [\texttt{arXiv:2209.14734}].

Contrastive learning has emerged as a dominant unsupervised paradigm for node-level tasks. Deep Graph Contrastive Representation Learning introduces a framework that maximizes the agreement between node representations derived from two corrupted graph views, targeting both structure and attribute levels [\texttt{arXiv:2006.04131}]. This method relies on mutual information and classical triplet loss to achieve robustness [\texttt{arXiv:2006.04131}]. To improve upon static corruption schemes, Graph Contrastive Learning with Adaptive Augmentation introduces a more nuanced mechanism; it employs node centrality measures for topological changes and noise injection for attribute corruption [\texttt{arXiv:2010.14945}]. By using importance-based feature corruption, this adaptive approach better preserves intrinsic structures and semantic information compared to standard augmentation techniques [\texttt{arXiv:2010.14945}].

Beyond static node representations, learning can also address temporal dynamics. The dyngraph2vec approach utilizes a deep architecture of dense and recurrent layers to learn temporal transitions [\texttt{arXiv:1809.02657}]. This enables the prediction of unseen links in dynamic networks by capturing the evolution of graph structures [\texttt{arXiv:1809.02657}]. Collectively, these advancements demonstrate a transition from simple structural reconstruction to sophisticated models capable of handling temporal evolution, adaptive feature corruption, and large-scale discrete diffusion.

\subsection{Explainability and Interpretability}

The interpretability of Graph Neural Networks (GNNs) has evolved from identifying structural importance to developing complex taxonomies of explanation methodologies. Early efforts focused on identifying the specific components that drive model predictions. GNNExplainer provides a model-agnostic approach for any graph-based machine learning task by maximizing the mutual information between a GNN's prediction and the distribution of possible subgraph structures; this allows for the identification of a compact subgraph and a subset of crucial node features via an optimization task \texttt{arXiv:1903.03894}. In contrast, GraphLIME offers a local interpretable framework specifically for node-level tasks by utilizing a nonlinear feature selection mechanism. Rather than focusing on subgraph topology, GraphLIME employs the Hilbert-Schmidt Independence Criterion (HSIC) Lasso on an \$N\$-hop neighborhood to identify the \$K\$ most representative features, making it applicable across diverse domains such as biology and social networks \texttt{arXiv:2001.06216}.

Further diversifying the landscape of model-agnostic explainers, PGM-Explainer shifts the focus from feature selection or subgraph extraction to probabilistic dependencies. By utilizing Probabilistic Graphical Models (PGM) and the concept of the Markov blanket, it generates explanations that demonstrate feature dependencies through conditional probabilities \texttt{arXiv:2010.05788}. While GNNExplainer and GraphLIME primarily isolate "what" is important (subgraphs or features), PGM-Explainer provides a more nuanced view of "how" features interact.

As the field has matured, the literature has shifted toward synthesizing these disparate methods into cohesive frameworks. Recent surveys have attempted to organize the growing body of work. One survey provides a comprehensive taxonomy and reviews performance metrics to standardize how explainability is evaluated \texttt{arXiv:2207.12599}. Another recent survey proposes a novel hierarchy that categorizes techniques based on both their underlying objectives and their specific methodologies, covering applications ranging from drug discovery to recommendation systems \texttt{arXiv:2306.01958}.

However, the pursuit of explainability must be balanced against the rigor of model evaluation. It is critical to note that the perceived superiority of complex GNN architectures—which these explainers aim to interpret—can sometimes be an artifact of flawed evaluation strategies. Research has demonstrated that improper data splits and training procedures can lead to misleading conclusions, often showing that simpler architectures can outperform more sophisticated ones when subjected to fair tuning \texttt{arXiv:1811.05868}. This suggests that interpretability efforts must be grounded in robust evaluation frameworks to ensure that the explanations provided are truly reflective of a model's learned logic rather than artifacts of an unsound training regime.

\subsection{The Frontier: Graph Transformers and LLM Integration}

The evolution of graph representation learning has progressed from specialized structural modeling to the integration of Large Language Models (LLMs) for complex reasoning. A primary frontier in this evolution is the development of Graph Transformers (GTs), which aim to reconcile the global receptive field of Transformers with the inherent structural constraints of graphs. The Structure-Aware Transformer (SAT) represents a significant advancement in this area by integrating structural information directly into the self-attention mechanism through subgraph representations extracted via a base GNN, achieving state-of-the-art performance on five graph prediction benchmarks \texttt{arXiv:2202.03036}. While SAT focuses on graph-level tasks, NodeFormer addresses the scalability challenges of transformer-style networks for node classification. By employing a kernelized Gumbel-Softmax operator, NodeFormer achieves linear algorithmic complexity for all-pair message passing, allowing for scalable structure learning on graphs containing up to 2 million nodes \texttt{arXiv:2306.08385}. Recent surveys further categorize these advancements, noting that GT architectures rely on mechanisms such as graph tokenization, positional encoding, and structure-aware attention to process diverse data types including molecules, proteins, and traffic networks \texttt{arXiv:2502.16533}.

Parallel to the refinement of GTs is the emergence of LLM integration, which shifts the focus toward textual graph understanding and reasoning. The literature identifies three primary technical roles for LLMs in this context: acting as a predictor, an encoder, or an aligner \texttt{arXiv:2312.02783}. This integration is particularly vital for text-attributed and text-paired graphs, such as academic or e-commerce networks. A critical challenge in this domain is managing the tension between graph complexity and LLM context windows. G-Retriever addresses this through a Retrieval-Augmented Generation (RAG) approach, utilizing Prize-Collecting Steiner Tree optimization for subgraph retrieval and soft prompting for fine-tuning to mitigate hallucinations \texttt{arXiv:2402.07630}.

While G-Retriever focuses on subgraph retrieval for question answering, GraphRAG introduces a different paradigm for global sensemaking. Instead of localized retrieval, GraphRAG builds an entity knowledge graph and generates pre-generated community summaries, enabling the processing of large datasets up to 1 million tokens \texttt{arXiv:2404.16130}. These specialized implementations are unified by a broader movement toward holistic GraphRAG frameworks that integrate query processors, retrievers, organizers, and generators \texttt{arXiv:2501.00309}. Ultimately, the synergy between these fields suggests a trajectory where graphs serve as structured knowledge sources that enhance LLM capabilities through advanced prompting (e.g., Chain-of-Thought or Graph-of-Thought) and structured data understanding across domains like relational databases and code \texttt{arXiv:2605.02452}.

\subsection{Tensions and Open Problems}

A primary theoretical tension exists regarding the fundamental limits of graph representation learning, specifically concerning the relationship between architectural design and expressive power. While the seminal work in \texttt{arXiv:1810.00826} established a framework for measuring GNN power via the Weisfeiler-Lehman test, subsequent literature has questioned whether this metric captures the true utility of neural architectures. For instance, \texttt{arXiv:1901.00596} suggests that the focus on WL-equivalence may overlook broader practical capabilities, while \texttt{arXiv:2003.04078} provides a more nuanced survey of expressive power that attempts to reconcile these formal bounds with real-world performance. This debate extends into the structural domain, where researchers disagree on how to best incorporate spatial information; while \texttt{arXiv:2009.00142} advocates for provably powerful distance encodings to enhance expressivity, \texttt{arXiv:2110.07875} argues for the necessity of learnable structural and positional features to capture complex topological nuances that fixed encodings might miss.

A second significant conflict arises in the scalability and robustness of GNNs when applied to non-idealized graph structures. There is a growing divide between methods optimized for homophilous, smooth signals and those designed for the "messy" reality of large-scale, heterophilous networks. As noted in \texttt{arXiv:2305.06142}, standard feature expansion techniques often struggle to bridge the gap identified in the non-homophilous benchmarks presented in \texttt{arXiv:2110.14446}. This tension is further exacerbated by the "bottleneck" problem, where the ability of GNNs to capture long-range dependencies is contested. While \texttt{arXiv:1710.10903} relies on attention mechanisms to weigh neighbor importance, \texttt{arXiv:2206.08164} demonstrates through the Long Range Graph Benchmark that even sophisticated attention-based models often fail to overcome the over-smoothing and over-squashing issues highlighted in \texttt{arXiv:2006.05205}.

Finally, the field faces a methodological rift concerning the interpretability and explainability of graph models. As GNNs are increasingly deployed in sensitive domains, the lack of a unified explanatory framework has become a critical bottleneck. Recent surveys, such as \texttt{arXiv:2306.01958}, critique existing taxonomic approaches like those in \texttt{arXiv:2012.15445} and \texttt{arXiv:2207.12599} for failing to provide rigorous, actionable metrics for what constitutes a "true" explanation. This lack of consensus on explainability is compounded by the difficulty of distinguishing between a model's learned structural logic and its reliance on spurious correlations, leaving the field in a state of flux between purely predictive success and verifiable transparency.

\subsubsection{Outlook}

The frontier of graph learning is rapidly shifting toward the integration of graph structures with large-scale generative intelligence. We are moving away from isolated message-passing paradigms toward hybrid architectures, such as Graph Retrieval-Augmented Generation (GraphRAG) as explored in \texttt{arXiv:2501.00309} and the synthesis of graph structures with LLMs in \texttt{arXiv:2506.02452}. The next generation of models will likely prioritize temporal continuity, as seen in the development of efficient continuous prediction models like Coden (\texttt{arXiv:2602.12613}), and structural diversity through parameter-free message passing (\texttt{arXiv:2508.19884}). Ultimately, the field is converging on a unified "Graph-Language" era, where the topological rigor of GNNs meets the reasoning capabilities of foundation models to solve complex, multi-modal reasoning tasks.

\section*{References}
\begin{enumerate}\small
  \item The Emerging Field of Signal Processing on Graphs: Extending
  High-Dimensional Data Analysis to Networks and Other Irregular Domains (2012). \texttt{arXiv:1211.0053}
  \item Spectral Networks and Locally Connected Networks on Graphs (2013). \texttt{arXiv:1312.6203}
  \item Adam: A Method for Stochastic Optimization (2014). \texttt{arXiv:1412.6980}
  \item DeepWalk: Online Learning of Social Representations (2014). \texttt{arXiv:1403.6652}
  \item Learning Phrase Representations using RNN Encoder-Decoder for
  Statistical Machine Translation (2014). \texttt{arXiv:1406.1078}
  \item Gated Graph Sequence Neural Networks (2015). \texttt{arXiv:1511.05493}
  \item Batch Normalization: Accelerating Deep Network Training by Reducing
  Internal Covariate Shift (2015). \texttt{arXiv:1502.03167}
  \item Convolutional Networks on Graphs for Learning Molecular Fingerprints (2015). \texttt{arXiv:1509.09292}
  \item LINE: Large-scale Information Network Embedding (2015). \texttt{arXiv:1503.03578}
  \item Deep Residual Learning for Image Recognition (2015). \texttt{arXiv:1512.03385}
  \item Deep Convolutional Networks on Graph-Structured Data (2015). \texttt{arXiv:1506.05163}
  \item Diffusion-Convolutional Neural Networks (2015). \texttt{arXiv:1511.02136}
  \item Semi-Supervised Classification with Graph Convolutional Networks (2016). \texttt{arXiv:1609.02907}
  \item Convolutional Neural Networks on Graphs with Fast Localized Spectral
  Filtering (2016). \texttt{arXiv:1606.09375}
  \item node2vec: Scalable Feature Learning for Networks (2016). \texttt{arXiv:1607.00653}
  \item Learning Convolutional Neural Networks for Graphs (2016). \texttt{arXiv:1605.05273}
  \item Revisiting Semi-Supervised Learning with Graph Embeddings (2016). \texttt{arXiv:1603.08861}
  \item Variational Graph Auto-Encoders (2016). \texttt{arXiv:1611.07308}
  \item Group Equivariant Convolutional Networks (2016). \texttt{arXiv:1602.07576}
  \item Geometric deep learning on graphs and manifolds using mixture model CNNs (2016). \texttt{arXiv:1611.08402}
  \item Geometric deep learning: going beyond Euclidean data (2016). \texttt{arXiv:1611.08097}
  \item Interaction Networks for Learning about Objects, Relations and Physics (2016). \texttt{arXiv:1612.00222}
  \item Categorical Reparameterization with Gumbel-Softmax (2016). \texttt{arXiv:1611.01144}
  \item Inductive Representation Learning on Large Graphs (2017). \texttt{arXiv:1706.02216}
  \item Graph Attention Networks (2017). \texttt{arXiv:1710.10903}
  \item Neural Message Passing for Quantum Chemistry (2017). \texttt{arXiv:1704.01212}
  \item Attention Is All You Need (2017). \texttt{arXiv:1706.03762}
  \item Modeling Relational Data with Graph Convolutional Networks (2017). \texttt{arXiv:1703.06103}
  \item Representation Learning on Graphs: Methods and Applications (2017). \texttt{arXiv:1709.05584}
  \item Dynamic Edge-Conditioned Filters in Convolutional Neural Networks on
  Graphs (2017). \texttt{arXiv:1704.02901}
  \item Stochastic Training of Graph Convolutional Networks with Variance
  Reduction (2017). \texttt{arXiv:1710.10568}
  \item MoleculeNet: A Benchmark for Molecular Machine Learning (2017). \texttt{arXiv:1703.00564}
  \item Diffusion Convolutional Recurrent Neural Network: Data-Driven Traffic
  Forecasting (2017). \texttt{arXiv:1707.01926}
  \item Deep Gaussian Embedding of Graphs: Unsupervised Inductive Learning via
  Ranking (2017). \texttt{arXiv:1707.03815}
  \item Graph Convolutional Matrix Completion (2017). \texttt{arXiv:1706.02263}
  \item CayleyNets: Graph Convolutional Neural Networks with Complex Rational
  Spectral Filters (2017). \texttt{arXiv:1705.07664}
  \item Residual Gated Graph ConvNets (2017). \texttt{arXiv:1711.07553}
  \item How Powerful are Graph Neural Networks? (2018). \texttt{arXiv:1810.00826}
  \item Hierarchical Graph Representation Learning with Differentiable Pooling (2018). \texttt{arXiv:1806.08804}
  \item On the Generalization of Equivariance and Convolution in Neural Networks
  to the Action of Compact Groups (2018). \texttt{arXiv:1802.03690}
  \item Tensor field networks: Rotation- and translation-equivariant neural
  networks for 3D point clouds (2018). \texttt{arXiv:1802.08219}
  \item Representation Learning on Graphs with Jumping Knowledge Networks (2018). \texttt{arXiv:1806.03536}
  \item Graph Convolutional Neural Networks for Web-Scale Recommender Systems (2018). \texttt{arXiv:1806.01973}
  \item Graph Neural Networks: A Review of Methods and Applications (2018). \texttt{arXiv:1812.08434}
  \item Relational inductive biases, deep learning, and graph networks (2018). \texttt{arXiv:1806.01261}
  \item Pitfalls of Graph Neural Network Evaluation (2018). \texttt{arXiv:1811.05868}
  \item Weisfeiler and Leman Go Neural: Higher-order Graph Neural Networks (2018). \texttt{arXiv:1810.02244}
  \item Deeper Insights into Graph Convolutional Networks for Semi-Supervised
  Learning (2018). \texttt{arXiv:1801.07606}
  \item Deep Graph Infomax (2018). \texttt{arXiv:1809.10341}
  \item 3D Steerable CNNs: Learning Rotationally Equivariant Features in
  Volumetric Data (2018). \texttt{arXiv:1807.02547}
  \item GraphRNN: Generating Realistic Graphs with Deep Auto-regressive Models (2018). \texttt{arXiv:1802.08773}
  \item FastGCN: Fast Learning with Graph Convolutional Networks via Importance
  Sampling (2018). \texttt{arXiv:1801.10247}
  \item Predict then Propagate: Graph Neural Networks meet Personalized PageRank (2018). \texttt{arXiv:1810.05997}
  \item GraphVAE: Towards Generation of Small Graphs Using Variational
  Autoencoders (2018). \texttt{arXiv:1802.03480}
  \item Invariant and Equivariant Graph Networks (2018). \texttt{arXiv:1812.09902}
  \item Learning Deep Generative Models of Graphs (2018). \texttt{arXiv:1803.03324}
  \item Intertwiners between Induced Representations (with Applications to the
  Theory of Equivariant Neural Networks) (2018). \texttt{arXiv:1803.10743}
  \item Link Prediction Based on Graph Neural Networks (2018). \texttt{arXiv:1802.09691}
  \item GaAN: Gated Attention Networks for Learning on Large and Spatiotemporal
  Graphs (2018). \texttt{arXiv:1803.07294}
  \item A General Theory of Equivariant CNNs on Homogeneous Spaces (2018). \texttt{arXiv:1811.02017}
  \item DynGEM: Deep Embedding Method for Dynamic Graphs (2018). \texttt{arXiv:1805.11273}
  \item Dynamic Graph CNN for Learning on Point Clouds (2018). \texttt{arXiv:1801.07829}
  \item Large-Scale Learnable Graph Convolutional Networks (2018). \texttt{arXiv:1808.03965}
  \item Attention-based Graph Neural Network for Semi-supervised Learning (2018). \texttt{arXiv:1803.03735}
  \item Adaptive Sampling Towards Fast Graph Representation Learning (2018). \texttt{arXiv:1809.05343}
  \item Deep Learning on Graphs: A Survey (2018). \texttt{arXiv:1812.04202}
  \item Adversarially Regularized Graph Autoencoder for Graph Embedding (2018). \texttt{arXiv:1802.04407}
  \item dyngraph2vec: Capturing Network Dynamics using Dynamic Graph
  Representation Learning (2018). \texttt{arXiv:1809.02657}
  \item Adaptive Graph Convolutional Neural Networks (2018). \texttt{arXiv:1801.03226}
  \item Streaming Graph Neural Networks (2018). \texttt{arXiv:1810.10627}
  \item Attention Models in Graphs: A Survey (2018). \texttt{arXiv:1807.07984}
  \item N-GCN: Multi-scale Graph Convolution for Semi-supervised Node
  Classification (2018). \texttt{arXiv:1802.08888}
  \item DynamicGEM: A Library for Dynamic Graph Embedding Methods (2018). \texttt{arXiv:1811.10734}
  \item GeniePath: Graph Neural Networks with Adaptive Receptive Paths (2018). \texttt{arXiv:1802.00910}
  \item A Comprehensive Survey on Graph Neural Networks (2019). \texttt{arXiv:1901.00596}
  \item Fast Graph Representation Learning with PyTorch Geometric (2019). \texttt{arXiv:1903.02428}
  \item PyTorch: An Imperative Style, High-Performance Deep Learning Library (2019). \texttt{arXiv:1912.01703}
  \item Simplifying Graph Convolutional Networks (2019). \texttt{arXiv:1902.07153}
  \item GNNExplainer: Generating Explanations for Graph Neural Networks (2019). \texttt{arXiv:1903.03894}
  \item Strategies for Pre-training Graph Neural Networks (2019). \texttt{arXiv:1905.12265}
  \item InfoGraph: Unsupervised and Semi-supervised Graph-Level Representation
  Learning via Mutual Information Maximization (2019). \texttt{arXiv:1908.01000}
  \item General \$E(2)\$-Equivariant Steerable CNNs (2019). \texttt{arXiv:1911.08251}
  \item Provably Powerful Graph Networks (2019). \texttt{arXiv:1905.11136}
  \item DropEdge: Towards Deep Graph Convolutional Networks on Node
  Classification (2019). \texttt{arXiv:1907.10903}
  \item Representation Learning for Dynamic Graphs: A Survey (2019). \texttt{arXiv:1905.11485}
  \item Deep Graph Library: A Graph-Centric, Highly-Performant Package for Graph
  Neural Networks (2019). \texttt{arXiv:1909.01315}
  \item Gauge Equivariant Convolutional Networks and the Icosahedral CNN (2019). \texttt{arXiv:1902.04615}
  \item GraphSAINT: Graph Sampling Based Inductive Learning Method (2019). \texttt{arXiv:1907.04931}
  \item Cluster-GCN: An Efficient Algorithm for Training Deep and Large Graph
  Convolutional Networks (2019). \texttt{arXiv:1905.07953}
  \item A Fair Comparison of Graph Neural Networks for Graph Classification (2019). \texttt{arXiv:1912.09893}
  \item Graph Neural Networks Exponentially Lose Expressive Power for Node
  Classification (2019). \texttt{arXiv:1905.10947}
  \item Heterogeneous Graph Attention Network (2019). \texttt{arXiv:1903.07293}
  \item EvolveGCN: Evolving Graph Convolutional Networks for Dynamic Graphs (2019). \texttt{arXiv:1902.10191}
  \item On the Universality of Invariant Networks (2019). \texttt{arXiv:1901.09342}
  \item Graph Neural Networks for Social Recommendation (2019). \texttt{arXiv:1902.07243}
  \item Self-Attention Graph Pooling (2019). \texttt{arXiv:1904.08082}
  \item Efficient Graph Generation with Graph Recurrent Attention Networks (2019). \texttt{arXiv:1910.00760}
  \item Relational Pooling for Graph Representations (2019). \texttt{arXiv:1903.02541}
  \item Position-aware Graph Neural Networks (2019). \texttt{arXiv:1906.04817}
  \item Universal Invariant and Equivariant Graph Neural Networks (2019). \texttt{arXiv:1905.04943}
  \item Explainability Techniques for Graph Convolutional Networks (2019). \texttt{arXiv:1905.13686}
  \item MixHop: Higher-Order Graph Convolutional Architectures via Sparsified
  Neighborhood Mixing (2019). \texttt{arXiv:1905.00067}
  \item DeepGCNs: Can GCNs Go as Deep as CNNs? (2019). \texttt{arXiv:1904.03751}
  \item Measuring and Relieving the Over-smoothing Problem for Graph Neural
  Networks from the Topological View (2019). \texttt{arXiv:1909.03211}
  \item Multi-scale Attributed Node Embedding (2019). \texttt{arXiv:1909.13021}
  \item On the equivalence between graph isomorphism testing and function
  approximation with GNNs (2019). \texttt{arXiv:1905.12560}
  \item Diffusion Improves Graph Learning (2019). \texttt{arXiv:1911.05485}
  \item Revisiting Graph Neural Networks: All We Have is Low-Pass Filters (2019). \texttt{arXiv:1905.09550}
  \item What graph neural networks cannot learn: depth vs width (2019). \texttt{arXiv:1907.03199}
  \item Graph U-Nets (2019). \texttt{arXiv:1905.05178}
  \item Graph Transformer Networks (2019). \texttt{arXiv:1911.06455}
  \item LanczosNet: Multi-Scale Deep Graph Convolutional Networks (2019). \texttt{arXiv:1901.01484}
  \item B-Spline CNNs on Lie Groups (2019). \texttt{arXiv:1909.12057}
  \item PairNorm: Tackling Oversmoothing in GNNs (2019). \texttt{arXiv:1909.12223}
  \item Weisfeiler and Leman go sparse: Towards scalable higher-order graph
  embeddings (2019). \texttt{arXiv:1904.01543}
  \item Approximation Ratios of Graph Neural Networks for Combinatorial Problems (2019). \texttt{arXiv:1905.10261}
  \item Understanding Attention and Generalization in Graph Neural Networks (2019). \texttt{arXiv:1905.02850}
  \item Node Embedding over Temporal Graphs (2019). \texttt{arXiv:1903.08889}
  \item Are Powerful Graph Neural Nets Necessary? A Dissection on Graph
  Classification (2019). \texttt{arXiv:1905.04579}
  \item Learning Discrete Structures for Graph Neural Networks (2019). \texttt{arXiv:1903.11960}
  \item Graphs, Convolutions, and Neural Networks: From Graph Filters to Graph
  Neural Networks (2020). \texttt{arXiv:2003.03777}
  \item Geom-GCN: Geometric Graph Convolutional Networks (2020). \texttt{arXiv:2002.05287}
  \item Graph Neural Networks: Architectures, Stability and Transferability (2020). \texttt{arXiv:2008.01767}
  \item Architectural Implications of Graph Neural Networks (2020). \texttt{arXiv:2009.00804}
  \item Retrieval-Augmented Generation for Knowledge-Intensive NLP Tasks (2020). \texttt{arXiv:2005.11401}
  \item REALM: Retrieval-Augmented Language Model Pre-Training (2020). \texttt{arXiv:2002.08909}
  \item Open Graph Benchmark: Datasets for Machine Learning on Graphs (2020). \texttt{arXiv:2005.00687}
  \item Benchmarking Graph Neural Networks (2020). \texttt{arXiv:2003.00982}
  \item Graph Contrastive Learning with Augmentations (2020). \texttt{arXiv:2010.13902}
  \item Contrastive Multi-View Representation Learning on Graphs (2020). \texttt{arXiv:2006.05582}
  \item KILT: a Benchmark for Knowledge Intensive Language Tasks (2020). \texttt{arXiv:2009.02252}
  \item Temporal Graph Networks for Deep Learning on Dynamic Graphs (2020). \texttt{arXiv:2006.10637}
  \item Deep Graph Contrastive Representation Learning (2020). \texttt{arXiv:2006.04131}
  \item A Generalization of Transformer Networks to Graphs (2020). \texttt{arXiv:2012.09699}
  \item Parameterized Explainer for Graph Neural Network (2020). \texttt{arXiv:2011.04573}
  \item Inductive Representation Learning on Temporal Graphs (2020). \texttt{arXiv:2002.07962}
  \item Graph Contrastive Learning with Adaptive Augmentation (2020). \texttt{arXiv:2010.14945}
  \item GCC: Graph Contrastive Coding for Graph Neural Network Pre-Training (2020). \texttt{arXiv:2006.09963}
  \item Generalizing Convolutional Neural Networks for Equivariance to Lie
  Groups on Arbitrary Continuous Data (2020). \texttt{arXiv:2002.12880}
  \item TUDataset: A collection of benchmark datasets for learning with graphs (2020). \texttt{arXiv:2007.08663}
  \item XGNN: Towards Model-Level Explanations of Graph Neural Networks (2020). \texttt{arXiv:2006.02587}
  \item PGM-Explainer: Probabilistic Graphical Model Explanations for Graph
  Neural Networks (2020). \texttt{arXiv:2010.05788}
  \item Explainability in Graph Neural Networks: A Taxonomic Survey (2020). \texttt{arXiv:2012.15445}
  \item Beyond Homophily in Graph Neural Networks: Current Limitations and
  Effective Designs (2020). \texttt{arXiv:2006.11468}
  \item Principal Neighbourhood Aggregation for Graph Nets (2020). \texttt{arXiv:2004.05718}
  \item On the Bottleneck of Graph Neural Networks and its Practical
  Implications (2020). \texttt{arXiv:2006.05205}
  \item Simple and Deep Graph Convolutional Networks (2020). \texttt{arXiv:2007.02133}
  \item Distance Encoding: Design Provably More Powerful Neural Networks for
  Graph Representation Learning (2020). \texttt{arXiv:2009.00142}
  \item Improving Graph Neural Network Expressivity via Subgraph Isomorphism
  Counting (2020). \texttt{arXiv:2006.09252}
  \item Can Graph Neural Networks Count Substructures? (2020). \texttt{arXiv:2002.04025}
  \item GraphLIME: Local Interpretable Model Explanations for Graph Neural
  Networks (2020). \texttt{arXiv:2001.06216}
  \item GPT-GNN: Generative Pre-Training of Graph Neural Networks (2020). \texttt{arXiv:2006.15437}
  \item Graph-Bert: Only Attention is Needed for Learning Graph Representations (2020). \texttt{arXiv:2001.05140}
  \item Adaptive Universal Generalized PageRank Graph Neural Network (2020). \texttt{arXiv:2006.07988}
  \item Interpreting Graph Neural Networks for NLP With Differentiable Edge
  Masking (2020). \texttt{arXiv:2010.00577}
  \item SIGN: Scalable Inception Graph Neural Networks (2020). \texttt{arXiv:2004.11198}
  \item Higher-Order Explanations of Graph Neural Networks via Relevant Walks (2020). \texttt{arXiv:2006.03589}
  \item A Survey on The Expressive Power of Graph Neural Networks (2020). \texttt{arXiv:2003.04078}
  \item Generalization and Representational Limits of Graph Neural Networks (2020). \texttt{arXiv:2002.06157}
  \item Permutation Invariant Graph Generation via Score-Based Generative
  Modeling (2020). \texttt{arXiv:2003.00638}
  \item Wiki-CS: A Wikipedia-Based Benchmark for Graph Neural Networks (2020). \texttt{arXiv:2007.02901}
  \item APAN: Asynchronous Propagation Attention Network for Real-time Temporal
  Graph Embedding (2020). \texttt{arXiv:2011.11545}
  \item Random Features Strengthen Graph Neural Networks (2020). \texttt{arXiv:2002.03155}
  \item Machine Learning on Graphs: A Model and Comprehensive Taxonomy (2020). \texttt{arXiv:2005.03675}
  \item Lorentz Group Equivariant Neural Network for Particle Physics (2020). \texttt{arXiv:2006.04780}
  \item Self-supervised Learning on Graphs: Deep Insights and New Direction (2020). \texttt{arXiv:2006.10141}
  \item DeeperGCN: All You Need to Train Deeper GCNs (2020). \texttt{arXiv:2006.07739}
  \item Heterogeneous Graph Transformer (2020). \texttt{arXiv:2003.01332}
  \item When Does Self-Supervision Help Graph Convolutional Networks? (2020). \texttt{arXiv:2006.09136}
  \item Building powerful and equivariant graph neural networks with structural
  message-passing (2020). \texttt{arXiv:2006.15107}
  \item Self-supervised Learning: Generative or Contrastive (2020). \texttt{arXiv:2006.08218}
  \item The Surprising Power of Graph Neural Networks with Random Node
  Initialization (2020). \texttt{arXiv:2010.01179}
  \item LieTransformer: Equivariant self-attention for Lie Groups (2020). \texttt{arXiv:2012.10885}
  \item Towards Deeper Graph Neural Networks (2020). \texttt{arXiv:2007.09296}
  \item Data Augmentation for Graph Neural Networks (2020). \texttt{arXiv:2006.06830}
  \item Foundations and modelling of dynamic networks using Dynamic Graph Neural
  Networks: A survey (2020). \texttt{arXiv:2005.07496}
  \item Graph Neural Networks: Methods, Applications, and Opportunities (2021). \texttt{arXiv:2108.10733}
  \item Ego-GNNs: Exploiting Ego Structures in Graph Neural Networks (2021). \texttt{arXiv:2107.10957}
  \item Graph Transformer Networks: Learning Meta-path Graphs to Improve GNNs (2021). \texttt{arXiv:2106.06218}
  \item Learning Graph Representations (2021). \texttt{arXiv:2102.02026}
  \item E(n) Equivariant Graph Neural Networks (2021). \texttt{arXiv:2102.09844}
  \item Geometric Deep Learning: Grids, Groups, Graphs, Geodesics, and Gauges (2021). \texttt{arXiv:2104.13478}
  \item Do Transformers Really Perform Bad for Graph Representation? (2021). \texttt{arXiv:2106.05234}
  \item Rethinking Graph Transformers with Spectral Attention (2021). \texttt{arXiv:2106.03893}
  \item On Explainability of Graph Neural Networks via Subgraph Explorations (2021). \texttt{arXiv:2102.05152}
  \item Geometric and Physical Quantities Improve E(3) Equivariant Message
  Passing (2021). \texttt{arXiv:2110.02905}
  \item Graph Contrastive Learning Automated (2021). \texttt{arXiv:2106.07594}
  \item A Practical Method for Constructing Equivariant Multilayer Perceptrons
  for Arbitrary Matrix Groups (2021). \texttt{arXiv:2104.09459}
  \item Adversarial Graph Augmentation to Improve Graph Contrastive Learning (2021). \texttt{arXiv:2106.05819}
  \item Graph Self-Supervised Learning: A Survey (2021). \texttt{arXiv:2103.00111}
  \item Large-Scale Representation Learning on Graphs via Bootstrapping (2021). \texttt{arXiv:2102.06514}
  \item Understanding over-squashing and bottlenecks on graphs via curvature (2021). \texttt{arXiv:2111.14522}
  \item Graph Neural Networks with Learnable Structural and Positional
  Representations (2021). \texttt{arXiv:2110.07875}
  \item How Attentive are Graph Attention Networks? (2021). \texttt{arXiv:2105.14491}
  \item Self-Supervised Learning of Graph Neural Networks: A Unified Review (2021). \texttt{arXiv:2102.10757}
  \item Beyond Low-frequency Information in Graph Convolutional Networks (2021). \texttt{arXiv:2101.00797}
  \item Large Scale Learning on Non-Homophilous Graphs: New Benchmarks and
  Strong Simple Methods (2021). \texttt{arXiv:2110.14446}
  \item GraphiT: Encoding Graph Structure in Transformers (2021). \texttt{arXiv:2106.05667}
  \item Frame Averaging for Invariant and Equivariant Network Design (2021). \texttt{arXiv:2110.03336}
  \item Geometric Deep Learning and Equivariant Neural Networks (2021). \texttt{arXiv:2105.13926}
  \item Self-supervised Learning on Graphs: Contrastive, Generative,or
  Predictive (2021). \texttt{arXiv:2105.07342}
  \item Scalars are universal: Equivariant machine learning, structured like
  classical physics (2021). \texttt{arXiv:2106.06610}
  \item Graph-level Neural Networks: Current Progress and Future Directions (2022). \texttt{arXiv:2205.15555}
  \item A Survey of Explainable Graph Neural Networks: Taxonomy and Evaluation
  Metrics (2022). \texttt{arXiv:2207.12599}
  \item Explainability in Graph Neural Networks: An Experimental Survey (2022). \texttt{arXiv:2203.09258}
  \item Geometrically Equivariant Graph Neural Networks: A Survey (2022). \texttt{arXiv:2202.07230}
  \item GreaseLM: Graph REASoning Enhanced Language Models for Question
  Answering (2022). \texttt{arXiv:2201.08860}
  \item Recipe for a General, Powerful, Scalable Graph Transformer (2022). \texttt{arXiv:2205.12454}
  \item DiGress: Discrete Denoising diffusion for graph generation (2022). \texttt{arXiv:2209.14734}
  \item Structure-Aware Transformer for Graph Representation Learning (2022). \texttt{arXiv:2202.03036}
  \item ROLAND: Graph Learning Framework for Dynamic Graphs (2022). \texttt{arXiv:2208.07239}
  \item Representing Long-Range Context for Graph Neural Networks with Global
  Attention (2022). \texttt{arXiv:2201.08821}
  \item Towards Better Evaluation for Dynamic Link Prediction (2022). \texttt{arXiv:2207.10128}
  \item SimGRACE: A Simple Framework for Graph Contrastive Learning without Data
  Augmentation (2022). \texttt{arXiv:2202.03104}
  \item Equiformer: Equivariant Graph Attention Transformer for 3D Atomistic
  Graphs (2022). \texttt{arXiv:2206.11990}
  \item TGL: A General Framework for Temporal GNN Training on Billion-Scale
  Graphs (2022). \texttt{arXiv:2203.14883}
  \item Evaluating Explainability for Graph Neural Networks (2022). \texttt{arXiv:2208.09339}
  \item SPECTRE: Spectral Conditioning Helps to Overcome the Expressivity Limits
  of One-shot Graph Generators (2022). \texttt{arXiv:2204.01613}
  \item Long Range Graph Benchmark (2022). \texttt{arXiv:2206.08164}
  \item Pure Transformers are Powerful Graph Learners (2022). \texttt{arXiv:2207.02505}
  \item Cooperative Graph Neural Networks (2023). \texttt{arXiv:2310.01267}
  \item A Survey on Explainability of Graph Neural Networks (2023). \texttt{arXiv:2306.01958}
  \item Feature Expansion for Graph Neural Networks (2023). \texttt{arXiv:2305.06142}
  \item Graph Rewriting for Graph Neural Networks (2023). \texttt{arXiv:2305.18632}
  \item A Survey on Spectral Graph Neural Networks (2023). \texttt{arXiv:2302.05631}
  \item Think-on-Graph: Deep and Responsible Reasoning of Large Language Model
  on Knowledge Graph (2023). \texttt{arXiv:2307.07697}
  \item Reasoning on Graphs: Faithful and Interpretable Large Language Model
  Reasoning (2023). \texttt{arXiv:2310.01061}
  \item Knowledge-Augmented Language Model Prompting for Zero-Shot Knowledge
  Graph Question Answering (2023). \texttt{arXiv:2306.04136}
  \item Retrieval-Augmented Generation for Large Language Models: A Survey (2023). \texttt{arXiv:2312.10997}
  \item REPLUG: Retrieval-Augmented Black-Box Language Models (2023). \texttt{arXiv:2301.12652}
  \item Large Language Models on Graphs: A Comprehensive Survey (2023). \texttt{arXiv:2312.02783}
  \item Chain-of-Knowledge: Grounding Large Language Models via Dynamic
  Knowledge Adapting over Heterogeneous Sources (2023). \texttt{arXiv:2305.13269}
  \item Knowledge Graph Prompting for Multi-Document Question Answering (2023). \texttt{arXiv:2308.11730}
  \item Towards Better Dynamic Graph Learning: New Architecture and Unified
  Library (2023). \texttt{arXiv:2303.13047}
  \item Temporal Graph Benchmark for Machine Learning on Temporal Graphs (2023). \texttt{arXiv:2307.01026}
  \item Enhancing Retrieval-Augmented Large Language Models with Iterative
  Retrieval-Generation Synergy (2023). \texttt{arXiv:2305.15294}
  \item NodeFormer: A Scalable Graph Structure Learning Transformer for Node
  Classification (2023). \texttt{arXiv:2306.08385}
  \item One for All: Towards Training One Graph Model for All Classification
  Tasks (2023). \texttt{arXiv:2310.00149}
  \item Retrieve-Rewrite-Answer: A KG-to-Text Enhanced LLMs Framework for
  Knowledge Graph Question Answering (2023). \texttt{arXiv:2309.11206}
  \item A Survey of Geometric Graph Neural Networks: Data Structures, Models and Applications (2024). \texttt{arXiv:2403.00485}
  \item Graphs Unveiled: Graph Neural Networks and Graph Generation (2024). \texttt{arXiv:2403.13849}
  \item Foundations and Frontiers of Graph Learning Theory (2024). \texttt{arXiv:2407.03125}
  \item DGNN: Decoupled Graph Neural Networks with Structural Consistency
  between Attribute and Graph Embedding Representations (2024). \texttt{arXiv:2401.15584}
  \item Scalable Graph Compressed Convolutions (2024). \texttt{arXiv:2407.18480}
  \item G-Retriever: Retrieval-Augmented Generation for Textual Graph
  Understanding and Question Answering (2024). \texttt{arXiv:2402.07630}
  \item From Local to Global: A Graph RAG Approach to Query-Focused
  Summarization (2024). \texttt{arXiv:2404.16130}
  \item Graph Retrieval-Augmented Generation: A Survey (2024). \texttt{arXiv:2408.08921}
  \item GNN-RAG: Graph Neural Retrieval for Large Language Model Reasoning (2024). \texttt{arXiv:2405.20139}
  \item KG-Agent: An Efficient Autonomous Agent Framework for Complex Reasoning
  over Knowledge Graph (2024). \texttt{arXiv:2402.11163}
  \item EdgeGFL: Rethinking Edge Information in Graph Feature Preference
  Learning (2025). \texttt{arXiv:2502.02302}
  \item Parameter-Free Structural-Diversity Message Passing for Graph Neural Networks (2025). \texttt{arXiv:2508.19884}
  \item Graph Neural Networks Powered by Encoder Embedding for Improved Node Learning (2025). \texttt{arXiv:2507.11732}
  \item A Survey of Graph Transformers: Architectures, Theories and Applications (2025). \texttt{arXiv:2502.16533}
  \item Graph Neural Networks for Databases: A Survey (2025). \texttt{arXiv:2502.12908}
  \item Retrieval-Augmented Generation with Graphs (GraphRAG) (2025). \texttt{arXiv:2501.00309}
  \item Invariants in equivariant birational geometry (2026). \texttt{arXiv:2602.23998}
  \item GraphNNK -- Graph Classification and Interpretability (2026). \texttt{arXiv:2602.00753}
  \item On Languages Describing Large Graph Classes (2026). \texttt{arXiv:2604.19719}
  \item Coden: Efficient Temporal Graph Neural Networks for Continuous Prediction (2026). \texttt{arXiv:2602.12613}
  \item Position: How can Graphs Help Large Language Models? (2026). \texttt{arXiv:2605.02452}
\end{enumerate}

\end{document}
